% =====================================================================
\chapter{Szakirodalmi áttekintés}
\label{ch:szakirodalom}
% =====================================================================

Ebből a fejezetből az olvasó megtudja, mit vár el a szakirodalom egy jó feleletválasztós kérdéstől, mely
megközelítések próbálták eddig gépi úton előállítani, és ezek mit oldanak meg a bevezetés négy
problémájából: a kulcs helyességéből (P1), a vizsga szerkezetéből és lefedettségéből (P2), az
adatvédelemből (P3) és a nagy, távoli modellektől való függésből (P4). A munkákat nem egyenként
ismertetjük, hanem azt vizsgáljuk, mit visznek közelebb a megoldáshoz, és hol hagynak rést. A módszertani
irodalom után a magyar nyelvtechnológiai előzmények, a jogi keret (GDPR, Infotv., MI-rendelet) és a
kutatási rés összevetése következik.

% ---------------------------------------------------------------------
\section{A jó feleletválasztós kérdés és a Bloom-taxonómia}
\label{sec:irod-mcq}
% ---------------------------------------------------------------------

A feleletválasztós kérdés (\emph{multiple-choice question}, MCQ) három részből áll: a \emph{törzsből}
(a kérdés vagy a befejezendő állítás), a \emph{kulcsból} (a helyes válasz) és a
\emph{disztraktorokból} (a helytelen, de hihető válaszlehetőségek). Népszerűségét az adja, hogy gépileg,
gyorsan és objektíven javítható; ára az, hogy jó kérdést írni nehéz.

A kérdésírás szabályait Haladyna, Downing és Rodriguez~\cite{haladyna2002review} foglalta össze 31
irányelvben. A tartalmi irányelvek szerint egy kérdés egyetlen, fontos tanulási célt mérjen; a formaiak
szerint a törzs önmagában is érthető legyen és kerülje a tagadást; a disztraktorokra vonatkozók szerint
a helytelen opciók legyenek hihetők, egymást kizárók, és formailag ne különbözzenek feltűnően a
kulcstól, nehogy „árulkodó jelet” adjanak. Hogy ezek nem esztétikai kérdések, azt Downing~\cite{downing2005effects}
egy orvosi vizsgán mutatta meg: a szabálysértő kérdések miatt a hallgatók 10--15\,\%-a került a
„nem felelt meg” kategóriába, akik a hibás kérdések nélkül átmentek volna. Egy gépi rendszernél tehát
nem elég, hogy a kérdések „jónak tűnjenek”: a szabályok betartását ellenőrizni és mérni kell. Ez adja a
P1 tartalmát: a helyes kérdéshez igaz kulcs, hamis, de hihető disztraktorok és árulkodó jel nélküli törzs
kell.

A disztraktorok számáról Rodriguez~\cite{rodriguez2005three} metaanalízise azt találta, hogy a legtöbb
esetben három opció is elegendő, mert a negyedik és ötödik ritkán „működik”. Tarrant és
munkatársai~\cite{tarrant2009distractors} szerint a disztraktorok jelentős része \emph{nem funkcionáló},
vagyis a hallgatók kevesebb mint 5\,\%-a választja, Gierl és munkatársai~\cite{gierl2017distractors}
áttekintése pedig rámutat, hogy az ilyen disztraktor gyakorlatilag csökkenti az opciók számát, és
növeli a találgatással elérhető pontszámot. A \Mimir{} a Moodle-ban megszokott négyopciós formát
követi, de a disztraktorok \emph{minőségét} -- hamisak és témába vágók-e -- külön mutatóval méri
(\ref{sec:mutatok}.~szakasz).

A kérdések kognitív szintjének leírására Bloom taxonómiáját használjuk Anderson és Krathwohl
átdolgozott változatában~\cite{anderson2001bloom}. Ez hat, egymásra épülő szintet különít el:
\emph{emlékezés}, \emph{megértés}, \emph{alkalmazás}, \emph{elemzés}, \emph{értékelés} és
\emph{alkotás}; feleletválasztós kérdéssel az első öt mérhető. Scully~\cite{scully2017constructing} szerint
magasabb rendű gondolkodás is mérhető így, ha a törzs új helyzetet ír le, a gyakorlatban mégis a
tényfelidéző kérdések dominálnak, mert ezeket a legkönnyebb megírni. Ugyanez várható egy nyelvi
modelltől is, ha maga dönti el, milyen kérdést ír. A \Mimir{} ezért a kért nehézséghez előre rögzített
Bloom-szint-eloszlást rendel, és azt a vizsgatervben kérdésenként kiosztja (\ref{sec:blueprint}.~szakasz);
a lefedettség (P2) így a témák mellett a kognitív szintekre is vonatkozik.

Hogy a kérdésírás milyen munkaigényes, azt a hazai tudásmérés is mutatja. A Szegedi Tudományegyetem eDia
online diagnosztikus mérési rendszerének~\cite{csapo2019edia} feladatbankjai összesen több mint 25\,000
feladatot tartalmaznak; ezeket képzett feladatírók készítik tartalmi keretrendszerek alapján, több
lépcsőben lektorálják, majd hallgatói mintán kipróbálják. Ehhez képest egy gépi rendszer \emph{vázlatot}
adhat: a lektorálást az oktatói jóváhagyás pótolhatja, a kipróbálást nem. A klasszikus
\emph{tételelemzés} (nehézség, megkülönböztető erő) ugyanis valódi hallgatói válaszokat igényel; a
generált vizsgák ilyen vizsgálata a további munka része (\ref{ch:osszefoglalas}.~fejezet).

% ---------------------------------------------------------------------
\section{Automatikus kérdésgenerálás: a szabályoktól a nagy nyelvi modellekig}
\label{sec:irod-aqg}
\label{sec:irod-llm-aqg}
% ---------------------------------------------------------------------

Az automatikus kérdésgenerálás (\emph{automatic question generation}, AQG) több évtizedes terület;
Kurdi és munkatársai~\cite{kurdi2020systematic} szisztematikus áttekintése 93 oktatási célú rendszert
elemzett. A fejlődésben három generáció különíthető el, és mindegyik másképp viszonyul a P1--P4
problémákhoz.

\paragraph{Szabály- és sablonalapú rendszerek.} Az első generáció a szöveget nyelvtanilag elemzi,
megkeresi a „kérdezhető” mondatrészeket, és transzformációs szabályokkal kérdő mondattá alakítja őket.
Mitkov és Ha rendszere~\cite{mitkov2003computer} a fontos főnévi csoportokat választja kulcsnak, és egy
lexikai adatbázis rokon fogalmait ajánlja disztraktornak. Heilman és Smith~\cite{heilman2010good} a
\emph{túlgenerálás és rangsorolás} elvét vezette be: a szabályok sok jelöltet adnak, amelyek közül egy
statisztikai modell választja ki a legjobbakat. Hazai példa a Miskolci Egyetem
rendszere~\cite{alshboul2024hybrid}, amely Python-programrészletekből ontológia és egy nyílt
kérdésgeneráló keretrendszer segítségével ír kérdéseket. Ezek a rendszerek helyben futnak, és a kulcs a
forrásból származik, így a P1-et részben, a P3-at és a P4-et teljesen megoldják. A kérdések viszont
merevek, többnyire mondat-átfogalmazások, a vizsga szerkezetéről (P2) nem döntenek, és minden nyelvhez
külön szabályrendszer kell -- a toldalékoló magyarhoz különösen sok.

\paragraph{Neurális rendszerek.} A második generációt a nagy, kézzel annotált adatkészletek tették
lehetővé, például a Wikipédia-bekezdésekhez írt, több mint százezer kérdés--válasz párt tartalmazó
SQuAD~\cite{rajpurkar2016squad}.
Du és munkatársai~\cite{du2017learning} ezeken tanítottak kódoló--dekódoló hálózatot, amely a bekezdésből
és a kiválasztott válaszból szóról szóra állítja elő a kérdést. A kérdések folyékonyabbak, de a
tanítóadat stílusát követik, és semmi sem garantálja, hogy a forrásból megválaszolhatók: a P1 itt
romlik, mert a törzs és a kulcs kapcsolata már tanult valószínűségből adódik.

A \emph{disztraktorgenerálás} külön irány; Alhazmi és munkatársai~\cite{alhazmi2024distractor} újabb
áttekintése a tudásbázisra épülő, a korpuszalapú és a
generatív módszereket veti össze. Megállapításuk szerint a generatív módszerek adják a legtermészetesebb
disztraktorokat, de gyakori hibájuk, hogy egy disztraktor valójában szintén helyes, vagy nyilvánvalóan
abszurd. Pontosan ezt a két hibatípust vizsgálja disztraktoronként a \Mimir{} ellenőrzője és a bíró.

\paragraph{Nagy nyelvi modellek.} A harmadik generációban a kérdéseket egy nagy nyelvi modell írja
utasítás alapján. Wang és munkatársai~\cite{wang2022towards} GPT-3-mal generáltak tankönyvi kérdéseket,
amelyek egy része emberi értékelés szerint az emberek által írtakkal összemérhető volt.
Doughty és munkatársai~\cite{doughty2024comparative} programozási kurzusokon 651 GPT-4-es és 449 oktatói
kérdést vetettek össze hatszempontos rubrikával. A gépi kérdések 81,7\,\%-a minden szempontnak
megfelelt, 4,9\,\%-ukban viszont több helyes válasz is volt (az oktatóiaknál 1,1\,\%), 4,0\,\%-ukban pedig
egy nyilvánvalóan hibás disztraktor elárulta a kulcsot (0,9\,\%); ezek átnézés nélkül a vizsgába kerültek
volna. Bitew és munkatársai~\cite{bitew2023distractor} a promptba tett, már jóváhagyott példakérdésekkel
javították a disztraktorokat; ugyanez a csoport készítette az OpenStax-tankönyvekhez írt, szakértők által
ellenőrzött EduQG adatkészletet~\cite{hadifar2023eduqg}, amelynek angol dokumentumait mi is használjuk. A
kis modellek oldaláról Fawzi és munkatársai~\cite{fawzi2023small} kimutatták, hogy egy kis modell oktatási
szövegeken végzett további tanítással a nagy modellekhez mérhető kérdéseket írhat.

\paragraph{Mit oldanak meg, és mit nem?} A harmadik generáció a nyelvi minőséget megoldotta, a P1-et
azonban csak statisztikailag javítja: a hibák aránya kicsi, de nem nulla, és a hibás kérdést semmi sem
jelöli meg. A P2-t a munkák nem vizsgálják, mert egyenként kérnek kérdést, a vizsgát pedig ember
állítja össze. A P3-at és a P4-et a felhőben futó nagy modellek előidézik; Fawzi és munkatársai kis
modellje kivétel, de minden új tárgynál további tanítást igényel. A kész kérdéseket utólag, emberi
értékeléssel mérik; arra, hogy a generálás \emph{folyamata} hogyan biztosíthatja a helyességet és a
lefedettséget, kevés figyelem jut. A \Mimir{} ezen a ponton kíván hozzájárulni.

\paragraph{A kiértékelés problémája.} Mulla és Gharpure~\cite{mulla2023aqg} szerint a rendszerek jelentős
részét a gépi fordításból átvett, referencia-átfedésen alapuló metrikákkal (BLEU, ROUGE) mérik, amelyek
azt számolják, hány szó egyezik a generált és egy ember írta referenciakérdés között. Egy vizsgakérdés
azonban egyetlen egyező szó nélkül is lehet kiváló, hibás kulccsal pedig használhatatlan. Kurdi és
munkatársai~\cite{kurdi2020systematic} azt is kimutatták, hogy a humán értékelések jellemzően kis mintán
készülnek, és ritkán ellenőrzik a kulcs helyességét. Munkánkban ezért a helyességet a forrással szembeni
ítéletekkel mérjük, a referenciakérdéseket pedig csak segédmutatókhoz használjuk
(\ref{ch:kiertekeles}.~fejezet).

% ---------------------------------------------------------------------
\section{Nagy nyelvi modellek röviden}
\label{sec:irod-llm}
% ---------------------------------------------------------------------

A mai nyelvi modellek a transzformer-architektúrára~\cite{vaswani2017attention} épülnek, és a szöveget
\emph{tokenekre} (szavakra vagy szódarabokra) bontva dolgozzák fel. A \emph{kódoló} modellek (például a
BERT~\cite{devlin2019bert}) egy szöveg jelentését egy vektorba (\emph{beágyazásba}) sűrítik, és keresésre
használják őket; a \emph{dekódoló}, generatív modellek (például a GPT-család~\cite{brown2020gpt3}) a
szöveget tokenenként állítják elő, a csevegő változatokat pedig utasításkövetésre is
hangolják~\cite{ouyang2022instructgpt}. A \Mimir{} kódoló modellel darabol és keres, generatív modellel ír
kérdést. A toldalékoló magyar szavak több tokenre bomlanak, mint az angolok, így ugyanaz a tartalom
magyarul több helyet foglal a modell bemenetében (\emph{kontextusában}). Két szöveg hasonlóságát a
beágyazásaik koszinuszával mérjük:
\begin{equation}
\cos(\mathbf{u},\mathbf{v})=\frac{\mathbf{u}\cdot\mathbf{v}}{\lVert\mathbf{u}\rVert\,\lVert\mathbf{v}\rVert},
\label{eq:cos}
\end{equation}
amelynek értéke 1, ha a két vektor azonos irányú, és annál kisebb, minél jobban eltér a két szöveg
jelentése. A darabolás, a keresés, a duplikátumszűrés és a lefedettség mérése is erre épül.

\paragraph{Hallucináció.}\label{sec:irod-hall}
A nyelvi modellek legfontosabb gyengesége a \emph{hallucináció}: folyékony, magabiztos, de a forrással
vagy a valósággal ellentétes szöveg előállítása. Ji és munkatársai~\cite{ji2023hallucination} két típust
különítenek el. Az \emph{intrinzikus} hallucináció ellentmond a forrásnak (például a forrás szerint a
kávé a 17.~században jutott el Európába, a modell a 9.~századot írja), az \emph{extrinzikus} pedig olyan
állítás, amely a forrásból sem nem igazolható, sem nem cáfolható. A
kérdésgenerálásban mindkét típus súlyos: ha a kulcs ellentmond a forrásnak, a felkészült hallgató veszít
pontot; ha a forrásban nem szerepel, a kérdés a tananyagból nem válaszolható meg. Mivel a hallucináció a
valószínűségi szövegelőállítás következménye, nem alkalmi hiba, a kulcs helyességét a forrásból kell
eldönteni, nem a modell magabiztosságából. Ezért tekintjük a legfontosabb mutatónak az
\emph{alátámasztottságot}, vagyis azt, hogy a forrás igazolja-e a kulcsot.

\paragraph{Strukturált kimenet.} A generátornak gépileg feldolgozható JSON-kimenetet kell adnia, de egy
kis modell időnként hibás JSON-t ír. Ezt vagy utólagos ellenőrzéssel és újrakéréssel, vagy
\emph{korlátozott dekódolással} lehet kezelni, amely generálás közben csak a formátumnak megfelelő
tokeneket engedi meg~\cite{willard2023guided}. A \Mimir{} az előbbit alkalmazza: a formai hibát az
ellenőrző első lépése szűri.

% ---------------------------------------------------------------------
\section{Visszakereséssel kiegészített generálás}
\label{sec:irod-rag}
% ---------------------------------------------------------------------

A hallucináció elleni első stratégia, hogy a modell ne csak a paramétereiben tárolt tudásra
támaszkodjon. A visszakereséssel kiegészített generálás (\emph{retrieval-augmented generation}, RAG)
során a rendszer a válasz előtt releváns szövegrészleteket keres egy gyűjteményből, és ezeket a
promptba téve kéri a választ. A RAG nevet Lewis és munkatársai~\cite{lewis2020rag} vezették be, akik egy sűrű keresőt és egy generatív
modellt kapcsoltak össze.

Gao és munkatársai~\cite{gao2023ragsurvey} áttekintése három paradigmát különít el. A \emph{naiv RAG}
egyetlen keresést és egyetlen generálást végez: a lekérdezéshez legjobban hasonló néhány részletet a
promptba másolja. A \emph{fejlett RAG} a keresés előtt és után is javít (például átírja a lekérdezést vagy
újrarangsorolja a találatokat). A \emph{moduláris RAG} a keresési, generálási és ellenőrzési lépéseket
szabadon kombinálja. A \Mimir{} tervezett generálási folyamata ebben a szóhasználatban moduláris RAG: a
keresés a vizsgaterv minden pontjára külön fut, és az ellenőrzés eredménye visszahat rá.

\paragraph{Lexikális és sűrű keresés.} A \emph{lexikális} keresés szavak alapján indexel; legismertebb
képviselője a BM25~\cite{robertson2009bm25}, amely azokat a részleteket rangsorolja előre, amelyekben a
lekérdezés szavai gyakoriak, de a gyűjteményben ritkák. Gyors, nem igényel tanítást, és jól kezeli a
szakkifejezéseket és tulajdonneveket, de a szinonimákat és a magyar ragozott alakokat nem ismeri fel. A
\emph{sűrű} keresés a szövegeket egy közös vektortérbe képezi, és a hasonlóságot \eqref{eq:cos} szerinti
koszinusszal méri. A DPR~\cite{karpukhin2020dpr} megmutatta, hogy ez a nyílt kérdésmegválaszolásban
felülmúlhatja a BM25-öt. Az általunk használt többnyelvű E5 modell~\cite{wang2024me5} nagy
szövegpárgyűjteményen, kontrasztív tanítással készült, és a magyart is támogatja. A keresők összehasonlító
mérései szerint azonban egyik módszer sem a legjobb minden adatkészleten~\cite{thakur2021beir}, és magyarra
sem egyértelmű a sűrű keresés előnye: a MILQA kérdés--válasz adatbázison a BM25 jobb visszakeresést adott
a vizsgált vektoros megoldásoknál~\cite{novak2023milqa}. Ezért gyakori a \emph{hibrid} keresés,
amely a két rangsort a reciprok rangfúzióval (RRF)~\cite{cormack2009rrf} egyesíti: egy részlet annál több
pontot kap, minél előrébb áll az egyes rangsorokban. A \Mimir{} egyik konfigurációja (\arm{E2h}) ezt a
változatot vizsgálja (\ref{sec:karok}.~szakasz).

\paragraph{Lekérdezés-átírás.} A keresés minősége a lekérdezésen is múlik: ha az rövid, általános vagy más
stílusú, mint a keresett szöveg, a legjobb kereső is rossz találatot ad. A HyDE~\cite{gao2023hyde} ezért
egy nyelvi modellel előbb egy „hipotetikus” választ írat, és annak beágyazásával keres, mert az jobban
hasonlít a keresett szövegre, mint maga a kérdés. Vizsgagenerálásban ez kulcskérdés. A naiv RAG lekérdezése ugyanis egy
\emph{feladatleírás} („10 feleletválasztós kérdés, könnyű”), amely semmit sem mond a keresett tartalomról.
A \Mimir{} vizsgaterve a tervező által javasolt tartalmi fogalmakkal keres, ami lényegében
lekérdezés-átírás. Eredményeink szerint ez a különbség áll a naiv RAG leggyakoribb kudarca mögött
(\ref{ch:eredmenyek}.~fejezet).

\paragraph{Darabolás.} A keresés egysége a \emph{darab}: a dokumentumot rövidebb részletekre kell bontani.
A rögzített hosszúságú ablak mondatok közepén is vághat; a \emph{szemantikus} darabolás ott vág, ahol a
szomszédos mondatok beágyazásainak hasonlósága hirtelen lecsökken, vagyis valószínűleg témaváltás van. Qu,
Tu és Bao~\cite{qu2024chunking} szerint a szemantikus darabolás előnye adatkészletfüggő és gyakran kisebb
a vártnál.
Ezért a darabolási stratégiát külön konfigurációkkal (\arm{E5a}--\arm{E5c}) vizsgáljuk, ahelyett hogy eleve
jobbnak tekintenénk a szemantikus változatot (\ref{sec:darabolas}.~szakasz). Eredményeink szerint a
szemantikus darabolás rövid, tartalom nélküli darabokat is létrehozhat, amelyekhez éppen a feladatleírás
típusú lekérdezés áll a legközelebb.

\paragraph{Forráshivatkozás.} A visszakeresés nem garantálja, hogy a válasz a visszakeresett szövegen
alapul. Ezért egyre több rendszer kéri, hogy a modell \emph{hivatkozzon} a felhasznált részletre. Az ALCE
benchmark~\cite{gao2023alce} azt méri, hogy a hivatkozások valóban alátámasztják-e az állításokat, és még a
legjobb modelleknél is sok hiányos vagy pontatlan hivatkozást talált. Rashkin és
munkatársai~\cite{rashkin2023ais} szerint egy állítás akkor \emph{tulajdonítható} egy forrásnak, ha egy
átlagos olvasó a forrás alapján egyetértene vele. A \Mimir{} minden kérdéshez megköveteli a kulcsot
igazoló darabok azonosítóját, és a bíró is ehhez hasonló, olvasóközpontú definícióval dönt az
alátámasztottságról.

\paragraph{Gráfalapú keresés és hosszú kontextus.} A GraphRAG~\cite{edge2024graphrag} a gyűjteményből
nyelvi modellel fogalom- és kapcsolatgráfot épít, és ennek csoportjai alapján válaszol; a \Mimir{} ennek
egy könnyített, egyetlen dokumentumra korlátozott változatát vizsgálja (\arm{E3}, \ref{sec:graf}.~szakasz).
A RAG valódi alternatívája az, hogy a teljes dokumentumot egyetlen promptban adjuk át, ahogy az oktatók
a gyakorlatban teszik. A névleges kontextushossz viszont nem azonos a ténylegesen kihasznált
hosszal: Liu és munkatársai~\cite{liu2024lost} szerint a modellek a bemenet elején és végén lévő
információt jobban használják, mint a közepén lévőt, a RULER benchmark~\cite{hsieh2024ruler} szerint pedig
a teljesítmény a hossz növelésével folyamatosan romlik. Li és munkatársai~\cite{li2024ragvslc}
összevetésében elegendő erőforrás mellett a hosszú kontextus átlagosan jobb, a RAG viszont lényegesen
olcsóbb, és a lekérdezések több mint 60\,\%-án (elemzésük szerint 63\,\%-án) a két megközelítés pontosan
azonos előrejelzést ad. Ebből két várakozás következik: a teljes
dokumentumot kapó modell kérdései a dokumentum egyes részeire tömörülhetnek (\ref{h:lefed}.~hipotézis), és
a hosszú kontextus egy kis kártyán a memóriaigény miatt korlátos (\ref{sec:vram}.~szakasz). Ugyanakkor egy
erős modellnél a teljes dokumentumos prompt komoly versenytárs, ezért ezt is mérjük.

\paragraph{Mit old meg a visszakeresés, és mit nem?} A RAG a P4-et kezeli a legközvetlenebbül: rövid
kontextussal működik, ezért kis, helyi modell is használhatja, ami a P3-at is lehetővé teszi. A P1-et
csak közvetve segíti, mert az ALCE mérései szerint a promptba tett forrásból nem következik, hogy a modell
arra támaszkodik. A P2-t a naiv RAG kifejezetten rontja, mert tartalom nélküli lekérdezéssel a dokumentum
kis részét látja. A fejlett és moduláris változatok egy-egy hiányt pótolnak, de jellemzően nagy modellel és
kérdésmegválaszolásban értékelték őket, nem vizsgagenerálásban.

% ---------------------------------------------------------------------
\section{Feladatfelbontás és ellenőrzés}
\label{sec:irod-felbontas}
% ---------------------------------------------------------------------

A dolgozat alapgondolata, hogy egy összetett feladatot rövid, jól körülhatárolt lépésekre bontva egy
kisebb modell is megbízhatóbban teljesít. Wei és munkatársai~\cite{wei2022cot} szerint a modell
összetett feladatokban pontosabb, ha lépésenkénti gondolatmenetet ír (\emph{chain-of-thought}). A
\emph{least-to-most} módszer~\cite{zhou2023leasttomost} a problémát előbb részfeladatokra bontja, majd
ezeket egymás után oldja meg, a \emph{Plan-and-Solve}~\cite{wang2023planandsolve} pedig arra kéri a
modellt, hogy előbb tervet készítsen, és csak utána hajtsa végre -- ez a vizsgaterv gondolatának
közvetlen megfelelője. Ezek a módszerek a felbontást a promptra bízzák, így a terv betartása ugyanúgy a
modell utasításkövetésén múlik, mint a közvetlen megközelítésben. Wu, Terry és
Cai~\cite{wu2022aichains} \emph{MI-láncai} ezzel szemben a programban kötik össze a nyelvi modell egymásra
épülő hívásait. Felhasználói vizsgálatuk szerint a láncok nemcsak jobb eredményt adtak, hanem
\emph{átláthatóbbak} és jobban \emph{irányíthatók} is voltak, mert minden köztes lépés látható és
módosítható.

A \Mimir{} a második irányba tartozik: a szerkezeti döntéseket (hány kérdés, melyik témából, milyen
szinten) a program hozza meg, a nyelvi modell csak jól körülhatárolt részfeladatot végez (fogalmak
javaslása, egy kérdés megírása, egy ítélet). Így minden lépés rövid kontextussal dolgozik, ami egy kis
modellnek kedvez (P4); a szerkezet a program által kikényszerített tulajdonság (P2); minden köztes
eredmény visszakövethető; és a lépések egyenként cserélhetők, ami lehetővé teszi az egytényezős kísérleti
tervet. Az ár a több modellhívás és a hosszabb futási idő. A hivatkozott eredményeket többnyire nagy
modellekkel, érvelési feladatokon mérték; hogy a felbontás kis modellnél vizsgagenerálásban mennyit ér,
azt \az{\ref{ch:eredmenyek}}. és \az{\ref{ch:modellek}}.~fejezet méri.

\paragraph{Ellenőrzés és önjavítás.}\label{sec:irod-ellenorzes}
A hallucináció elleni második stratégia a kimenet utólagos ellenőrzése. A Self-Refine~\cite{madaan2023selfrefine}
a modellel értékelteti és javíttatja saját kimenetét, a Chain-of-Verification~\cite{dhuliawala2024cove}
ellenőrző kérdéseket tesz fel a saját válaszáról, és az eltérések alapján javít. A FActScore~\cite{min2023factscore} a szöveget atomi tényekre bontja, és mindegyiket
külön veti össze egy tudásforrással. Egy másik ág a \emph{természetes nyelvi következtetés}
(\emph{natural language inference}, NLI): annak eldöntése, hogy egy állítás következik-e egy szövegből,
ellentmond-e neki, vagy egyik sem. Az ilyen modellek az összefoglalók tényszerűségének ellenőrzésére is
alkalmasak~\cite{laban2022summac}. Az alátámasztottság
eldöntése pontosan ilyen feladat: következik-e a forrásból, hogy a megjelölt kulcs a helyes válasz?

Az önellenőrzésnek azonban korlátai vannak. Huang és munkatársai~\cite{huang2024selfcorrect} szerint
\emph{külső visszajelzés nélkül} az önjavítás érvelési feladatokban nem javít, sőt ronthat is. Snell és
munkatársai~\cite{snell2025scaling} ugyanakkor kimutatták, hogy a futás közben elköltött többletszámítás
(több minta, ellenőrzés, újrapróbálkozás) bizonyos feladatokban helyettesítheti a nagyobb modellt. A két
eredmény együtt azt sugallja, hogy az ellenőrzés akkor hasznos, ha \emph{külső} információt visz a
folyamatba: determinisztikus szabályokat, a forrásszöveget vagy egy másik modell ítéletét. A \Mimir{}
ellenőrzője ezért olcsó szabályokat (forma, hivatkozás, szivárgás, duplikátum) kombinál a forrással
szembeni, nyelvi modellel végzett ítéletekkel, és a hibát visszajelzésként adja vissza a generátornak
(\ref{sec:verifier}.~szakasz).

A visszacsatolásnak saját veszélye is van. Goodhart törvénye Strathern megfogalmazásában: ha egy
mérőszám céllá válik, megszűnik jó mérőszám lenni~\cite{strathern1997}. Ha a generátort addig javíttatjuk,
amíg át nem megy az ellenőrzésen, olyan megoldást is találhat, amely az ellenőrzést kielégíti, a valódi
célt viszont nem. A pilot vizsgálatunk egy ilyen esetet tárt fel: a legolcsóbb út a forrásmondat szinte
szó szerinti bemásolása volt a törzsbe (\emph{szivárgás}, \ref{sec:szivargas}.~szakasz).

% ---------------------------------------------------------------------
\section{Kis nyelvi modellek helyi futtatása}
\label{sec:irod-kis}
% ---------------------------------------------------------------------

A helyi futtatás feltétele, hogy a modell elférjen egy fogyasztói grafikus kártya memóriájában. Ma
több nyílt modellcsalád ér el jó teljesítményt néhány milliárd paraméterrel, például a
Qwen2.5~\cite{qwen25}, a Llama~3.1~\cite{grattafiori2024llama3} és a Gemma~3~\cite{gemma3}; ezek a kisebb
méretet több és jobb minőségű tanítóadattal ellensúlyozzák. A pilotban és az előzetes vizsgálatban a
Qwen2.5-7B, a fő vizsgálatban a Gemma~3 12B és a NYTK magyarra továbbtanított
PULI-LlumiX 8B modellje~\cite{pulillumix} a generátor (a magyar modellekről
\az{\ref{sec:irod-magyar}}.~szakasz szól). Természetes várakozás, hogy egy magyarra hangolt modell magyar
tananyagból jobb kérdéseket ír, mint egy általános többnyelvű modell; ezt a fő vizsgálatban közvetlenül
teszteljük (\ref{ch:modellek}.~fejezet).

A memóriaigény fő összetevője a súlyok mérete. Egy 7 milliárd paraméteres modell 16 bites ábrázolásban
paraméterenként 2 bájtot, összesen 14\,GB-ot foglal. A \emph{kvantálás} a súlyokat kevesebb biten
tárolja: a GPTQ~\cite{frantar2023gptq} a tanítás után 3--4 bitre tömörít. A
helyi futtatás de facto szabványa a \texttt{llama.cpp}~\cite{gerganov2023llamacpp} GGUF formátuma és
csoportos K-kvantálása; az általunk használt Q4\_K\_M változat átlagosan kb.\ 4,8 bitet használ
súlyonként, és erre épül az Ollama futtatókörnyezet~\cite{ollama2024} is. Átfogó kiértékelések szerint a
4 bites súlykvantálás általános feladatokon csak kis pontosságvesztést okoz~\cite{li2024evalquant}. A
memóriaigény második összetevője a \emph{kulcs--érték (KV) gyorsítótár}, amely a már feldolgozott tokenek
köztes eredményeit tárolja, így a kontextus hosszával arányosan nő. Mindezek alapján egy 7--8B modell
4 bites változata egy 8\,GB-os, a Gemma~3 12B egy 12\,GB-os fogyasztói kártyán futtatható; a számszerű
becslést \az{\ref{sec:vram}}.~szakasz adja meg. A helyi futtatás tehát egyszerre kezeli a P3-at és a
P4-et, de csak akkor, ha a módszer a kis modell rövidebb kontextusához és gyengébb utasításkövetéséhez
igazodik; a kvantálási mérések pedig általános feladatokra szólnak, nem vizsgagenerálásra.

% ---------------------------------------------------------------------
\section{A kiértékelés: nyelvi modell mint bíró és a mérés megbízhatósága}
\label{sec:irod-biro}
% ---------------------------------------------------------------------

Az emberi értékelés drága és lassú, ezért elterjedt, hogy egy erős nyelvi modell pontozza a kimenetet
(\emph{LLM-as-a-judge}); a dolgozatban ezt a kérdéseket pontozó, független nyelvi modellt \emph{bírónak}
nevezzük. Zheng és munkatársai~\cite{zheng2023judge} szerint egy erős bíró ítéletei közel annyira egyeznek
az emberi preferenciákkal, mint az emberek ítéletei egymással. Chiang és Lee~\cite{chiang2023human} hasonló
eredményre jutott úgy, hogy a modell szó szerint ugyanazt az utasítást és értékelőlapot kapta, mint az
emberi szakértők -- ezt az elvet a mi rubrikánk is követi. A RAGAS~\cite{es2024ragas} a
RAG-rendszerek referencia nélküli, nyelvi modellel számolt hűség- és relevanciamutatóit javasolja. Ezek az
egyezési eredmények főként párosított preferenciaítéletekre (melyik válasz a jobb) szólnak; a mi bíránk
kérdésenként, egy forrásrészlethez mérve ítél, amire csak korlátozottan vihetők át.

A bírók torzításai jól dokumentáltak: a \emph{pozíciós torzítás} az első vagy második helyen álló válasz
előnyben részesítése~\cite{wang2024fair}, a \emph{hossztorzítás} a hosszabb válaszoké, az
\emph{önpreferencia} pedig a saját modellcsalád szövegeié~\cite{panickssery2024selfpref}. Ezért bíránk, a
gpt-oss-120b~\cite{gptoss2025} más modellcsaládba tartozik, mint bármelyik generátor
(\ref{sec:biro}.~szakasz), ami az önpreferenciát csökkenti.

A bíró csak akkor érvényes mérőeszköz, ha ítéletei egyeznek az emberi szakértőkével. Két értékelő
bináris ítéleteire Cohen-$\kappa$-t~\cite{cohen1960kappa}, több értékelőre és ordinális skálára
Krippendorff-$\alpha$-t~\cite{krippendorff2004} használnak; mindkettő azt méri, mennyivel nagyobb az
egyezés a véletlenül várhatónál. Landis és Koch~\cite{landis1977} skálája szerint $0{,}20$ alatt csekély,
$0{,}21$--$0{,}40$ között gyenge, $0{,}41$--$0{,}60$ között közepes, $0{,}60$ felett erős az egyezés;
Krippendorff szerint $\alpha\ge0{,}8$ megbízható, $0{,}667$ és $0{,}8$ között óvatos következtetésre
elegendő. A bíró oktatói validálását \az{\ref{sec:oktatoi}}.~szakasz írja le.

\paragraph{A kísérletek megbízhatósága.}\label{sec:irod-meres}
Bouthillier és munkatársai~\cite{bouthillier2021variance} kimutatták, hogy a gépi tanulási eredményekben a
véletlen tényezők okozta szórás gyakran összemérhető a közölt javulásokkal, így egyetlen futás alapján nem
szabad két módszer között dönteni. Dror és munkatársai~\cite{dror2018hitchhiker} szerint a
nyelvtechnológiai cikkek jelentős része nem végez szignifikanciatesztet, vagy az adatokhoz nem illő próbát
használ; a 0 és 1 közé szorított, ferde eloszlású mutatóknál a nemparaméteres próbák (például a
Wilcoxon-próba) és a bootstrap a biztonságosabbak. Fontos a mérés \emph{egysége} is: egy dokumentum tíz
kérdése nem tíz független megfigyelés, ezért mi a dokumentumot tekintjük egységnek, és a sok összevetésből
adódó véletlen „szignifikanciát” Holm-korrekcióval~\cite{holm1979} szűrjük (\ref{sec:statisztika}.~szakasz).
Veszélyt jelent az \emph{adatszennyezés} is~\cite{sainz2023contamination}: a nyilvános Wikipédia- és
OpenStax-szövegek a modellek tanítóadatában is szerepelhettek, így a bíró a kérdést a forrás nélkül,
emlékezetből is megválaszolhatja.

Végül a mérés az oktatót is érinti. Parasuraman és Manzey~\cite{parasuraman2010complacency} szerint a
felhasználók hajlamosak túlságosan megbízni az automatikus rendszerek jelzéseiben
(\emph{automatizálási torzítás}): ha egy rendszer egy kérdést „ellenőrzöttnek” jelöl, az oktató kevésbé
figyelmesen nézi át; a jelenséget az MI-rendelet is megnevezi (\ref{sec:irod-jog}.~szakasz). Ezért az
ellenőrző jelzéseit a \Mimir{} felületén figyelmeztetésként, nem ítéletként jelenítjük meg, amíg
megbízhatóságukat mérések nem igazolják.

% ---------------------------------------------------------------------
\section{A magyar nyelv és a hazai előzmények}
\label{sec:irod-magyar}
% ---------------------------------------------------------------------

A magyar két szempontból jelent külön kihívást. Agglutináló nyelv, így egy szótőnek sok alakja van
(„immunrendszer”, „immunrendszert”, „immunrendszerben”), ami a lexikális keresést nehezíti és a
tokenszámot növeli. Emellett kis erőforrású nyelv: a MEGA kiértékelés~\cite{ahuja2023mega} szerint a
generatív modellek teljesítménye a kisebb erőforrású nyelveken jelentősen elmarad az angoltól, ezért egy
kis modell magyar szövegében nyelvtani hibákra is számítani kell, amit a kiértékelésnek külön mérnie kell.

\paragraph{Magyar nyelvi modellek.} Az első széles körben használt magyar kódoló a
huBERT~\cite{nemeskey2021hubert} volt. A HUN-REN Nyelvtudományi Kutatóközpont (NYTK) PULI-családjának első
tagjairól az MSZNY 2023-on számoltak be~\cite{yang2023jonnek}: egy BERT-Large, egy GPT-2 és egy 6,7
milliárd paraméteres GPT-3 típusú modellről; a BERT-Large több feladaton felülmúlta a huBERT-et, a
GPT-3 típusú modellel pedig a szerzők az első magyar, néhány példás (\emph{few-shot}) promptolási
kísérleteket végezték. Ugyanebben az évben egy- és háromnyelvű GPT-modelljeiket is
közölték~\cite{yang2023puligpt3}. A dolgozatban használt PULI-LlumiX 8B a Llama~3.1 magyar szövegeken
folytatott további előtanításával készült~\cite{pulillumix}. A csoport 2025-ben a Llama~3.1 és a Qwen2.5
alapú, magyar utasításadatokon finomhangolt csevegő változatokat (ChatPULI) is leírt, amelyek mérésük
szerint a magyar benchmarkokon és a többfordulós párbeszédben jobbak a korábbi
PULI-modelleknél~\cite{yang2025chatpuli}. Ez a fő vizsgálat értelmezéséhez lényeges: mi az utólagos
utasítás-hangolás nélküli, továbbelőtanított változatot mértük, így a PULI-ra vonatkozó következtetésünk
erre a változatra szól (\ref{ch:modellek}.~fejezet).

\paragraph{Magyar benchmarkok és kérdés--válasz rendszerek.} A magyar nyelvértési feladatok egységes
mérésére a HuLU benchmark~\cite{ligetinagy2024hulu} szolgál. A MILQA~\cite{novak2023milqa,novak2023milqahu}
magyar Wikipédia-cikkekre épülő, az angol SQuAD~2.0 elveit követő kérdés--válasz adatbázis, eldöntendő,
többrészes és számolást igénylő kérdéstípusokkal bővítve; adatforrásként mi is használjuk
(\ref{sec:adatkeszlet}.~szakasz). Erre épül a szegedi SHunQA nyíltkérdés-megválaszoló
rendszer~\cite{berkecz2024shunqa}, amely a választ a forrásszövegben jelöli meg, hogy ellenőrizhető legyen.
Ezek a munkák kérdésmegválaszolással foglalkoznak, de két tanulságuk átvihető: magyar szövegen a lexikális
keresés erős alapvonal, és a forrásban megjelölt válasz teszi ellenőrizhetővé a rendszert.

\paragraph{Hazai kérdésgenerálás és oktatási MI-használat.} Magyar nyelvű vizsgakérdés-generálásról
publikált, rendszerszintű kiértékelést nem találtunk. A legközelebbi előzmény a Szegedi
Tudományegyetem 2024--2025-ös, HUN-REN Cloudon futó projektje, amely a helyben futtatható nyelvi modellek
tananyag alapú magyar kérdésgenerálását vizsgálta, és nehézségként a modellek eltérő magyar nyelvi
képességét és a kérdésekkel szembeni összetett elvárásokat nevezte meg~\cite{szte2025kerdesgen};
eredményeiről publikációt nem találtunk. Az igényt hazai felmérések is mutatják. Egy 1027 hallgatós
felmérésben a válaszadók kb.\ 87\,\%-a használt már MI-eszközt, az aktív használat akadályai között pedig
etikai (kb.\ 42\,\%) és adatvédelmi (kb.\ 30\,\%) aggályok is szerepeltek~\cite{rajki2024hallgatoi}. Egy
1103 egyetemi oktatót felölelő felmérés szerint az oktatók az MI-t főként rutinfeladatokra (fordítás,
nyelvi ellenőrzés) használják, a mélyebb pedagógiai felhasználás háttérbe szorul~\cite{tnagy2025oktatoi}.
A vizsgakészítés ilyen mélyebb felhasználás, amelyhez megbízható és adatvédelmi szempontból vállalható
eszköz kell.

% ---------------------------------------------------------------------
\section{Jogi keret: adatvédelem és MI-rendelet}
\label{sec:irod-jog}
% ---------------------------------------------------------------------

A P3 jogi kereteit két közvetlenül alkalmazandó uniós rendelet, a GDPR~\cite{gdpr2016} és az
MI-rendelet~\cite{aiact2024} adja, hazai kiegészítésekkel. Az áttekintés mérnöki olvasat, nem jogi
tanácsadás: azt vizsgálja, milyen tervezési követelmények következnek ezekből. A megvalósítás állapotát
\az{\ref{ch:modszer}}.~fejezet tárgyalja.

\subsection{Az általános adatvédelmi rendelet}

\paragraph{Szerepek.} A GDPR az \emph{adatkezelőt} (aki a célt és az eszközöket meghatározza) és az
\emph{adatfeldolgozót} (aki az adatkezelő nevében, utasítása szerint dolgozik) különíti el. Az Európai
Adatvédelmi Testület szerint a szerepet a tényleges tevékenység dönti el, és aki az adatot saját céljára is
felhasználja, abban a körben maga is adatkezelő~\cite{edpb2021controller}. Egy felhőben futó csevegő
asszisztens tehát akkor 28.~cikk szerinti adatfeldolgozó, ha az intézménnyel kötött szerződés utasításhoz,
titoktartáshoz, biztonsági intézkedésekhez, a további adatfeldolgozók jóváhagyásához és a szerződés végi
törléshez köti. Ha az oktató saját fogyasztói fiókkal tölti fel a jegyzetét, ilyen garancia nincs; a
szolgáltató címzett, és ha a tartalmat például modellfejlesztésre is felhasználja, e körben adatkezelő. A
tananyagban szereplő harmadik személyek adatai így tudtuk nélkül kerülnek új adatkezeléshez.

\paragraph{Alapelvek, beépített adatvédelem, biztonság.} Az 5.~cikk alapelvei közül az
\emph{adattakarékosság} (csak a szükséges adat) és a \emph{korlátozott tárolhatóság} (csak a szükséges
ideig) hat közvetlenül a tervezésre: a dokumentum szövegére csak a generálás idejére van szükség, ezért a
legkisebb kockázatú megoldás a feldolgozás utáni törlés. Az \emph{elszámoltathatóság} naplózást kíván, de
olyat, amely nem tárol tartalmat. A 25.~cikk \emph{beépített és alapértelmezett adatvédelmet} ír elő
(például a vizsga mentése csak kérésre történjen), a 32.~cikk a kockázattal arányos \emph{biztonsági
intézkedéseket} (álnevesítés, titkosítás, bizalmasság, integritás, rendszeres ellenőrzés). A helyi futtatás
a továbbítás kockázatát szünteti meg, a hozzáférés-szabályozást nem pótolja.

\paragraph{Hatásvizsgálat, továbbítás, automatizált döntés.} Valószínűleg magas kockázatú, különösen új
technológiát alkalmazó adatkezelés előtt \emph{adatvédelmi hatásvizsgálat} kell (35.~cikk); generatív MI
és esetleges hallgatói adatok mellett intézményi bevezetés előtt ez indokolt. A modell maga is kockázat:
a tanítóadat egyes részletei a kész modellből visszanyerhetők~\cite{carlini2021extracting}, és az EDPB
szerint egy MI-modell nem tekinthető automatikusan anonimnak~\cite{edpb2024aimodels}; ezért felhasználói
tartalommal nem tanítunk modellt. Az EGT-n kívülre személyes adat csak a V.~fejezet szerinti
megfelelőségi határozat vagy megfelelő garanciák mellett továbbítható. Az Egyesült Államokba irányuló
továbbítás alapja ma az EU--USA adatvédelmi keretrendszer~\cite{dpf2023}, amelynek érvényességéről
fellebbezési eljárás folyik~\cite{latombe2025appeal}; a helyi futtatás ettől független. Végül a 22.~cikk
korlátozza a kizárólag automatizált, joghatással járó döntéseket: a kérdésgenerálás ilyet nem hoz, a
hallgatói válaszok emberi felülvizsgálat nélküli pontozása viszont már ide tartozhat.

\subsection{Az Infotv. és a NAIH}

Az Infotv.\ (2011.~évi CXII.~törvény)~\cite{infotv2011} kijelöli a GDPR hatálya alá tartozó adatkezelésekre
alkalmazandó hazai kiegészítő szabályokat (2.~§ (2)~bekezdés), és létrehozza a felügyeleti hatóságot, a
NAIH-ot; hazai üzemeltetésnél neki kell az incidenseket bejelenteni. A NAIH figyelemfelhívása szerint
csevegőrobot használatakor lehetőleg ne töltsünk fel személyes adatot a szolgáltató felületére, mert ez
előre nem látható adatkezeléshez vezethet, amelyért a feltöltő személyesen is felelhet~\cite{naihfigyelem}.
A hazai bankok MI-használatáról szóló 2025-ös jelentése~\cite{naih2025bankok} az \emph{árnyék-MI}
jelenségét elemzi: a munkatársak a szervezet tudta nélkül használnak MI-eszközöket, így a szervezet
elveszítheti a kontrollt az általa kezelt adatok felett. Jó gyakorlatként a nem saját infrastruktúrán
futó eszközök korlátozását, illetve a felhasználási adatok modelltanítását kizáró vállalati szerződést
említi. A jegyzetét saját fiókkal nyilvános csevegő asszisztensbe feltöltő oktató is árnyék-MI-t
használ: a P3 a hazai hatóság által is megnevezett kockázat.

\subsection{Az MI-rendelet}

\paragraph{Hatály és kockázati osztályok.} Az MI-rendelet a forgalomba hozott vagy üzembe helyezett
rendszerekre vonatkozik; az azt megelőző kutatásra, tesztelésre és fejlesztésre -- a valós körülmények
közötti tesztelés kivételével -- nem (2.~cikk (8)~bekezdés). A \Mimir{} kiértékelő keretrendszere ebbe a
körbe tartozik, egy egyetemi üzembe helyezés viszont már nem; ekkor a fejlesztő \emph{szolgáltató}, az
intézmény \emph{alkalmazó}. A rendelet négy szintet különít el: (1) \emph{tiltott gyakorlatok} (5.~cikk),
például az érzelemfelismerés oktatási intézményben; (2) \emph{magas kockázatú} rendszerek (6.~cikk, I. és
III.~melléklet) kockázatkezelési, dokumentációs, naplózási és emberi felügyeleti követelményekkel; (3)
\emph{átláthatósági kötelezettségek} (50.~cikk); (4) \emph{minimális kockázat}, külön kötelezettség
nélkül. A beépített általános célú modellek (Qwen, Gemma, Llama-alapú PULI) kötelezettségei a modellek
szolgáltatóit terhelik.

\paragraph{Oktatás a III.~mellékletben.} A III.~melléklet 3.~pontja szerint magas kockázatú (a) a
felvétel eldöntése és a személyek intézményekhez rendelése; (b) a tanulási eredmények értékelése, akkor
is, ha az eredmény a tanulási folyamatot irányítja; (c) annak felmérése, milyen szintű oktatásban
részesülhet valaki; és (d) a tiltott vizsgaviselkedés megfigyelése. A 6.~cikk (3)~bekezdése szerint egy
ilyen rendszer mégsem magas kockázatú, ha nem jelent jelentős kockázatot, például mert egy értékelés
\emph{előkészítő feladatát} látja el -- kivéve, ha profilalkotást végez; a kivételre hivatkozó
értékelést dokumentálni kell (6.~cikk (4)~bekezdés).

\paragraph{Emberi felügyelet, MI-jártasság, átláthatóság.} A 14.~cikk szerint a magas kockázatú rendszer
felügyelője értse a rendszer képességeit és korlátait, legyen tudatában az automatizálási torzításnak, és
a kimenetet felülbírálhassa. A 4.~cikk eredetileg 2025.~február~2-tól a munkatársak megfelelő
\emph{MI-jártasságának} biztosítását írta elő; a 2026-os módosítás óta az MI-jártasság fejlesztésének
\emph{támogatását}, meghatározott szint nélkül~\cite{omnibus2026}. Az 50.~cikk szerint a felhasználót
tájékoztatni kell, ha MI-rendszerrel érintkezik (1.~bekezdés); a szintetikus szöveget generáló rendszer
szolgáltatója a kimenetet géppel olvasható módon MI-generáltként jelöli (2.~bekezdés); a nyilvánosság
tájékoztatására közzétett MI-szöveget pedig fel kell tüntetni, kivéve, ha emberi szerkesztői ellenőrzésen
ment át (4.~bekezdés).

\paragraph{Alkalmazási határidők.} A rendelet 2024.~augusztus~1-jén lépett hatályba, és lépcsőzetesen
alkalmazandó (\ref{tab:mi-hataridok}.~táblázat). A Bizottság 2025.~november~19-én „digitális omnibusz”
csomagban (COM(2025)~836) javasolta egyebek mellett a magas kockázatú kötelezettségek halasztását. Ezt
(EU)~2026/1744 rendeletként fogadták el; 2026.~július~24-én hirdették ki, és 2026.~július~27-én lépett
hatályba~\cite{omnibus2026}. A III.~melléklet szerinti magas kockázatú kötelezettségek így
2027.~december~2-tól alkalmazandók; az 50.~cikk átláthatósági kötelezettségeit viszont nem halasztották
el, a III.~melléklet 3.~pontját pedig a módosítás nem érintette. Mivel az ütemezés egyszer már változott,
éles bevezetés előtt a hatályos szöveget újra ellenőrizni kell. A hazai végrehajtásról a
2025.~évi LXXV.~törvény~\cite{mitv2025} rendelkezik: létrehozta a Magyar Mesterséges Intelligencia Tanácsot,
a NAIH jelentése szerint pedig a piacfelügyeleti feladatokat a nemzetgazdasági miniszter, a bejelentő
hatóságiakat a Nemzeti Akkreditáló Hatóság látja el~\cite{naih2025bankok}.

\begin{table}[tbp]
\centering
\caption[Az MI-rendelet alkalmazási időpontjai a 2026-os módosítás után.]{Az MI-rendelet alkalmazási
időpontjai a 2026-os módosítás (digitális omnibusz) után. Források:~\cite{aiact2024,omnibus2026}; a dőlt
sorokat a módosítás vezette be vagy változtatta meg.}
\label{tab:mi-hataridok}
\small
\begin{tabular}{@{}lp{0.68\textwidth}@{}}
\toprule
Időpont & Mi válik alkalmazandóvá\\
\midrule
2024.~aug.~1. & A rendelet hatálybalépése\\
2025.~febr.~2. & Általános rendelkezések, MI-jártasság (4.~cikk), tiltott gyakorlatok (5.~cikk)\\
2025.~aug.~2. & Az általános célú MI-modellek szolgáltatóinak kötelezettségei, irányítás, szankciók\\
2026.~júl.~27. & \emph{A módosító rendelet hatálybalépése, az új 4.~cikk}\\
2026.~aug.~2. & Általános alkalmazás, köztük az 50.~cikk átláthatósági kötelezettségei\\
2026.~dec.~2. & \emph{Az 50.~cikk (2) a korábban forgalomba hozott rendszerekre; új tilalmak (5.~cikk)}\\
2027.~dec.~2. & \emph{III.~melléklet szerinti magas kockázatú rendszerek (eredetileg 2026.~aug.~2.)}\\
2028.~aug.~2. & \emph{I.~melléklet szerinti magas kockázatú rendszerek (eredetileg 2027.~aug.~2.)}\\
\bottomrule
\end{tabular}
\end{table}

\subsection{A \Mimir{} besorolása}

A \Mimir{} kérdésgeneráló funkciója nem tartozik a III.~melléklet 3.~pontjába. Nem dönt felvételről vagy
képzési szintről, nem figyeli a hallgatókat, és nem értékel tanulási eredményt: az oktatónak ad vázlatot,
amelyet ő átnéz, szerkeszt, és csak kötelező jóváhagyás után exportálhat; a kimenet egy emberi döntés
alapanyaga, nem maga a döntés. Ha egy szigorú értelmezés a vizsga előkészítését az értékelés részének
tekintené, a 6.~cikk (3)~bekezdés d)~pontja (előkészítő feladat) akkor is kivételt adna, mert a rendszer
nem végez profilalkotást; ezt dokumentálni kell. A besorolás után is kötelező az 50.~cikk (a csevegőfelület
és az MI-generált kérdések jelölése), az MI-jártasság támogatása és a GDPR.

A besorolás a felhasználástól függ, nem a technikától. Ha a rendszer a hallgatói válaszokat automatikusan
pontozná, vagy a kérdések kiválasztásával a tanulási folyamatot irányítaná, az a 3.~b) pont szerint magas
kockázatú felhasználás lenne, a teljes követelményrendszerrel, és a GDPR 22.~cikke is felmerülne. Ezért
pontozási funkció nem kerül a rendszerbe, csak új kockázati értékelés után. A 14.~cikk logikáját a
jelenlegi besorolásban is követjük: az oktató a döntéshozó, az ellenőrző jelzései figyelmeztetések, mert az
automatizálási torzítást a rendelet és a szakirodalom~\cite{parasuraman2010complacency} egyaránt megnevezi;
az UNESCO útmutatója~\cite{unesco2023guidance} oktatáspolitikai szinten vár el emberi ellenőrzést és
adatvédelmet. A \Mimir{} tervezési céljai --
helyi futtatás, zéró adatmegőrzés, kötelező oktatói jóváhagyás -- ezekből következnek
(\ref{sec:architektura}.~szakasz).

% ---------------------------------------------------------------------
\section{Összegzés és a kutatási rés}
\label{sec:irod-res}
% ---------------------------------------------------------------------

\Az{\ref{tab:irodalom-osszevetes}}.~táblázat a tárgyalt megközelítéseket a bevezetés P1--P4 problémái
szerint veti össze. A korai rendszerek helyben futnak, de sem a helyességet, sem a szerkezetet nem
ellenőrzik; a csevegő modell egyik problémát sem kezeli; a naiv RAG forráshoz köti a generálást, de a
feladatleírásból készült lekérdezéssel a dokumentum kis részét látja; az önjavító, a hivatkozó és a
gráfalapú módszerek egy-egy problémát részben kezelnek, de nagy modellt feltételeznek. A \Mimir{} sorában a helyesség mérése
csak részben teljesül: a bíró oktatói validálása még hátravan, és a saját ellenőrzőjének kérdésenkénti
megbízhatósága is korlátozottnak bizonyult (\ref{ch:eredmenyek}.~fejezet).

\begin{table}[tbp]
\centering
\caption[A kérdésgenerálási megközelítések összevetése a P1--P4 problémák szempontjából.]{A kérdésgenerálási megközelítések összevetése a P1--P4 problémák szempontjából
(\checkmark: kezeli, $\circ$: részben, --: nem kezeli). A táblázat azt mutatja, hogy egy megközelítés
foglalkozik-e az adott szemponttal, nem azt, hogy milyen jól.}
\label{tab:irodalom-osszevetes}
\small
\setlength{\tabcolsep}{4pt}
\resizebox{\textwidth}{!}{%
\begin{tabular}{@{}lcccccc@{}}
\toprule
Megközelítés & \makecell{Kulcs\\ellenőrzése} & \makecell{Szerkezet,\\lefedettség} & \makecell{Helyi\\futás} & \makecell{Kis\\modell} & \makecell{Forrás-\\hivatkozás} & \makecell{Helyesség\\mérése}\\
\midrule
Szabály/sablon alapú~\cite{mitkov2003computer,heilman2010good} & $\circ$ & -- & \checkmark & \checkmark & $\circ$ & -- \\
Neurális seq2seq~\cite{du2017learning} & -- & -- & \checkmark & \checkmark & -- & -- \\
Dokumentum $\to$ csevegő LLM (status quo) & -- & -- & -- & -- & -- & -- \\
LLM-alapú kérdésgenerálás~\cite{wang2022towards,doughty2024comparative} & -- & $\circ$ & -- & -- & -- & $\circ$ \\
Naiv RAG~\cite{lewis2020rag} & -- & -- & $\circ$ & $\circ$ & $\circ$ & -- \\
Továbbtanított kis modell~\cite{fawzi2023small} & -- & -- & \checkmark & \checkmark & -- & $\circ$ \\
Feladatfelbontó promptolás~\cite{zhou2023leasttomost,wang2023planandsolve} & -- & $\circ$ & -- & -- & -- & -- \\
Önjavítás / CoVe~\cite{madaan2023selfrefine,dhuliawala2024cove} & $\circ$ & -- & -- & -- & -- & -- \\
Hivatkozó generálás~\cite{gao2023alce} & $\circ$ & -- & -- & -- & \checkmark & \checkmark \\
GraphRAG~\cite{edge2024graphrag} & -- & $\circ$ & -- & -- & $\circ$ & -- \\
\midrule
\textbf{\Mimir{} (jelen munka)} & \checkmark & \checkmark & \checkmark & \checkmark & \checkmark & $\circ$ \\
\bottomrule
\end{tabular}}
\end{table}

Az összevetésből három hiányosság rajzolódik ki. (i) Tudomásunk szerint nincs olyan nyílt rendszer, amely a
kérdésgenerálást explicit tervvel, kérdésenkénti forráshivatkozással és ellenőrzéssel, fogyasztói GPU-n
(8--12\,GB) futó kis modellel végezné. (ii) A kiértékelések ritkán mérik a kulcs helyességét és a vizsga
lefedettségét, és hiányoznak az egytényezős kísérleti tervek, amelyek a komponensek hatását egyenként
választják szét. (iii) A magyar nyelvű vizsgakérdés-generálás kiértékelése feltáratlan: a hazai
nyelvtechnológia erős modelleket és kérdés--válasz adatbázisokat épített, de publikált kérdésgenerálási
kiértékelést nem találtunk. Az első hiányosságra \az{\ref{ch:modszer}}., a másodikra
\az{\ref{ch:kiertekeles}}., a harmadikra -- a rendelkezésre álló adatok korlátai között --
\az{\ref{ch:eredmenyek}}. és \az{\ref{ch:modellek}}.~fejezet ad választ.

\paragraph{A fejezet tanulsága.} A szakirodalom mind a négy problémára kínál részmegoldást, de egyik
megközelítés sem kezeli egyszerre mindet. A nagy modellek a kérdések nyelvi minőségét megoldották, a kulcs
helyességét (P1) azonban csak statisztikailag javítják, a vizsga szerkezetét (P2) pedig nem vizsgálják. A
visszakeresés és a helyi futtatás a P3-at és a P4-et kezeli, de a naiv változat a P2-t rontja. A
feladatfelbontás és a forráshoz mért ellenőrzés a P1-et és a P2-t célozza, de eddig főként nagy
modellekkel vizsgálták. A jogi keret a helyi, emberi jóváhagyással működő megoldás felé mutat: a tanár
által ellenőrzött kérdésgenerálás nem magas kockázatú felhasználás, a hallgatói válaszok automatikus
pontozása az lenne.
